\section{Modelos matemáticos propuestos}

En esta primera fase se presentan los modelos que se pretenden utilizar para analizar el acceso a internet en los hogares del Perú.

\subsection{Modelo 1: acceso a internet total}

Primero se analizará el acceso total a internet en los hogares del Perú. Para ello se usará la variable $I(t)$, que representa el porcentaje total de hogares con internet. Como este porcentaje no puede crecer de forma infinita, se considerará un modelo que tenga como límite máximo el 100\%.

Para la comparación inicial se puede considerar el modelo exponencial:
\begin{equation}
    \frac{dI}{dt}=kI,
\end{equation}
donde $k$ representa la tasa de crecimiento del acceso total a internet.

Además, debido a que el porcentaje tiene un límite superior, se propone el modelo logístico:
\begin{equation}
    \frac{dI}{dt}=rI\left(1-\frac{I}{K}\right), \qquad K=100,
\end{equation}
donde $r$ es la tasa de crecimiento y $K$ representa la capacidad máxima expresada en porcentaje.

\subsection{Modelo 2: comparación entre zona urbana y zona rural}

Luego se estudiará el acceso a internet según área de residencia. Para esto se usarán las variables $U(t)$, que representa el porcentaje de hogares urbanos con internet, y $R(t)$, que representa el porcentaje de hogares rurales con internet. Esta comparación permitirá observar si el crecimiento del acceso a internet es similar o diferente entre ambas zonas.

El comportamiento de ambas áreas se puede representar mediante modelos logísticos independientes:
\begin{align}
    \frac{dU}{dt} &= r_U U\left(1-\frac{U}{K_U}\right), \\
    \frac{dR}{dt} &= r_R R\left(1-\frac{R}{K_R}\right).
\end{align}

\subsection{Metodología propuesta para la segunda fase}

En la segunda fase del trabajo se desarrollarán los cálculos matemáticos y la comparación de resultados. Los pasos serán:

\begin{itemize}
    \item Organizar los datos en función del tiempo $t$.
    \item Graficar los datos reales de internet total, internet urbano, internet rural y computadoras.
    \item Plantear el modelo de ecuación diferencial correspondiente.
    \item Estimar los parámetros de cada modelo.
    \item Comparar los resultados de los modelos con los datos reales.
    \item Interpretar los resultados obtenidos.
    \item Elaborar conclusiones.
\end{itemize}
